\subsection*{A.\ Participant Demographics}

All 177 participants live in the United States and had used at least one of the six agents, which were eligibility criteria. Table~\ref{tab:demographics} gives their age, gender, and education.

\begin{table}[t]
\centering
\caption{Participant demographics ($n=177$). Percentages are of the 177 participants.}
\label{tab:demographics}
\small
\begin{tabular}{@{}lrr@{}}
\toprule
 & $n$ & \% \\
\midrule
\multicolumn{3}{@{}l}{\textit{Age}} \\
18--24 years & 29 & 16.4 \\
25--34 years & 56 & 31.6 \\
35--44 years & 45 & 25.4 \\
45--54 years & 27 & 15.3 \\
55--64 years & 17 & 9.6 \\
65+ years & 3 & 1.7 \\
\midrule
\multicolumn{3}{@{}l}{\textit{Gender}} \\
Man & 104 & 58.8 \\
Woman & 72 & 40.7 \\
Prefer not to say & 1 & 0.6 \\
\midrule
\multicolumn{3}{@{}l}{\textit{Education}} \\
Bachelor's degree & 75 & 42.4 \\
Graduate or professional degree & 35 & 19.8 \\
Some college, but no degree & 30 & 16.9 \\
Associate's or technical degree & 19 & 10.7 \\
High school diploma or GED & 15 & 8.5 \\
Some high school or less & 3 & 1.7 \\
\bottomrule
\end{tabular}
\end{table}

\subsection*{B.\ Survey Instrument}

\subsubsection*{B.1.\ Scenarios}

Participants saw the two scenarios in random order, each followed by the same questions. We reproduce them word for word.

\begin{itemize}[leftmargin=*,nosep]
\item \textit{Less sensitive scenario.} Imagine that you are chatting with an AI chatbot and would like to know the weather of your location. The conversation may look like the example below: User: I am in Orlando now. What's the weather today in Orlando? AI: Based on your city, it's 72\,$^{\circ}$F and sunny. There is a slight chance of rain this evening. User: Should I bring an umbrella? AI: Probably not. The chance of rain is only 10\%. Now you want to delete your location data to protect your privacy.
\item \textit{More sensitive scenario.} Imagine that you are chatting with an AI chatbot and want to know what to prepare for medical center visit. The conversation might look like the example below: User: I live at 7482 Orange Blossom Lane, Orlando, FL 32801. Can you find me the nearest Medical Center to get treatment for sexually transmitted infections? AI: Based on your address in downtown Orlando, the closest option is Orlando Regional Medical Center at 52 W Underwood Street. If it's a life-threatening emergency, call 911 rather than driving yourself. User: How should I prepare? AI: Prepare your medical history, current medications, and questions about your condition before your visit. Now you want to delete your location/address data to protect your privacy.
\end{itemize}

\subsubsection*{B.2.\ Data Deletion Method Options}

Before the scenarios, participants saw a description of the methods, each with an example screenshot, and marked which they had used. The answer choices were: (1) Type a message asking the AI Chatbot to forget it. (2) Delete this single conversation. (3) Clear all conversation history. (4) Delete a specific saved memory or fact. (5) Clear all saved memories. (6) Delete my account entirely. (7) Privacy dashboard or account-level data management page. (8) File a request. (9) Other (Please specify). A further choice read ``I have never deleted my data in AI Chatbots.'' In the description above the question, the sixth method read ``Delete the account entirely.'' and the last two methods were both numbered 7.

In each scenario, the question ``Among the methods you have used, which would you use to delete your location data in the scenario above?'' offered only the methods the participant had marked as used. A later question, ``Among the data deletion methods, including those you have used before and those you have not, which would you choose to delete your location data in this scenario? (Select all that apply)'', offered all of them.

\subsubsection*{B.3.\ Scales}

In each scenario, participants rated the method they had chosen on three scales, using five-point agreement scales (Strongly Disagree, Disagree, Neutral, Agree, Strongly Agree). The first label of the effort scale was misspelled in the survey as ``Stongly Disagree''.

\begin{itemize}[leftmargin=*,nosep]
\item \textit{Perceived protection (5 items).} This data deletion method contributes to the protection of my privacy. This data deletion method helps me to protect my personal information. This data deletion method effectively reduces the risk of my data being misused. This data deletion method gives me complete control over my personal data. By using this data deletion method, I can decide what personal data the AI Chatbot retains about me.
\item \textit{Effort (4 items).} This data deletion method takes too much time. This data deletion method is overly complex to use. I find it difficult to navigate this data deletion method. Understanding how this data deletion method works requires significant effort.
\item \textit{Benefit loss (3 items).} Using this data deletion method would prevent AI Chatbots from providing responses tailored to my activity context. Using this data deletion method would reduce the relevance of assistance tailored to my preferences. Using this data deletion feature would prevent AI Chatbots from providing the kind of help I need.
\end{itemize}

\subsubsection*{B.4.\ Other Questions}

The survey ran in this order: consent; screening (age, U.S. residence, and prior use of at least one of the six agents); an agent usage profile (which agents, for how long) and the method-use question in the previous subsection; the two scenarios in random order; deletion experience; post-deletion feedback expectations; privacy awareness; social influence and attitudes; and demographics.

In each scenario, the questions were: the method choice; why they chose it; what they expected to happen to their data (select all that apply); whether they had verified it (select all that apply); their confidence in that answer; the three scales; their willingness to delete, from ``Definitely not'' to ``Definitely yes''; the message they would type to the agent and why they would phrase it that way; and the all-methods question. The expectation options were: (1) It is permanently deleted from the database and backups of the AI Chatbot. (2) It is made invisible to users, such as by hiding the conversation containing the data. (3) The AI chatbot will not be able to reference it (or its derivatives) in the current conversation. (4) The AI chatbot will not be able to reference it (or its derivatives) in future conversation. (5) It is no longer used for training the AI model. (6) I am not sure what happens. (7) Others (Please specify). The verification options were: (1) Yes, I asked the AI Chatbot in the same conversation and it no longer recalled my data. (2) Yes, I asked the AI Chatbot in a new conversation and it no longer recalled my data. (3) Yes, I checked the memory or saved data settings on the AI Chatbot. (4) Yes, I checked whether the conversation or data was still visible in the user interface. (5) Yes, I submitted a data access request through the AI Chatbot privacy portal. (6) No, I did not know how to check. (7) No, I don't want to check. (8) Others (Please specify).

The deletion-experience block asked what types of information participants typically share, which agents they had tried to delete data from, what data they tried to delete, why, how they found out about the deletion options, and whether they had noticed an agent referencing information they believed they had deleted. The feedback block asked what confirmation they receive after deleting data and what they think agents should provide. The privacy-awareness block asked whether they had read privacy policies and included a six-item knowledge quiz on the right to be forgotten, training data, and memory. The social-influence block asked how people in their social network feel about using agents and sharing data with them, how they feel about submitting personal information and using agents for tasks, and how well statements about their ability to work with AI describe them. Demographics covered education, gender, and primary domain of work or study.

\subsubsection*{B.5.\ Attention Checks}

Two instructed-response items checked attention: ``Please select `Sometimes' for this question.'' (options: Never, Rarely, Sometimes, Often, Always), asked in the deletion-experience block, and ``Please select `Agree' for this question.'' (options: Strongly Disagree to Strongly Agree), asked in the social-influence block.

\subsection*{C.\ Statistical Details}

\subsubsection*{C.1.\ Analyses, Tests, and Correction Families}

Table~\ref{tab:analysis-families} lists every analysis reported in the paper, the test, the responses it uses, and the family over which $p$ values are adjusted with the Holm correction. Where no family is given, the analysis is a single test and its $p$ value is not adjusted.

\begin{table*}[t]
\centering
\caption{Analyses, tests, and correction families. A response is one participant's answer in one scenario. GEE is a generalized estimating equation.}
\label{tab:analysis-families}
\small
\begin{tabular}{@{}p{0.34\textwidth}p{0.28\textwidth}p{0.18\textwidth}p{0.14\textwidth}@{}}
\toprule
Analysis & Test & Base & Holm family \\
\midrule
Willingness to delete, less vs more sensitive scenario & Wilcoxon signed-rank; rank-biserial $r$ & 177 pairs & none (single test) \\
Protection, effort, benefit loss, less vs more & Wilcoxon signed-rank; rank-biserial $r$ & 177 pairs & the 3 scales \\
Method chosen, less vs more (overall) & Stuart--Maxwell marginal homogeneity & 176 participants who chose a method in both scenarios & none (single test) \\
Method chosen, less vs more (each method) & McNemar & 177 participants & the 7 methods \\
Expectation options, less vs more & McNemar & 177 participants & the 7 options \\
Verification options, less vs more & McNemar & 177 participants & the 8 options \\
Protection, effort, benefit loss by method & Gaussian GEE, clustered by participant (Wald) & 317 responses, 4 methods & the 3 scales \\
Effort, pairwise between methods & GEE contrasts, adjusted for scenario & 317 responses, 4 methods & the 6 method pairs \\
Expectation options by method & Logistic GEE, clustered by participant & 317 responses, 4 methods & the 7 options \\
Verification options by method & Chi-square or Fisher's exact, each scenario separately & participants choosing a method with at least 14 participants in that scenario & the 8 options, within each scenario \\
Scale structure & Confirmatory factor analysis (Appendix~C.2) & 354 responses & not applicable \\
Scenario order & Binomial test, Mann--Whitney, Fisher's exact (Appendix~C.3) & 177 participants & the 4 ratings (Mann--Whitney) \\
\bottomrule
\end{tabular}
\end{table*}

\subsubsection*{C.2.\ Confirmatory Factor Analysis}

The three scales measure separate things, so we report each on its own. We fit a confirmatory factor analysis with three factors, one per scale (protection with 5 items, effort with 4, and benefit loss with 3), by maximum likelihood on the 354 responses. Standardized loadings ranged from .73 to .94, the average variance extracted was .64 to .76, and Cronbach's $\alpha$ was .87 to .92. On the comparative indices the fit was acceptable (CFI .94, TLI .92, SRMR .060), but the RMSEA, .099, is above the usual .08 cutoff, and the $\chi^2$ test is significant, as it commonly is with samples of this size. Because the pooled model treats the two responses of a participant as independent, we also fit each scenario alone ($n=177$ each), with similar results (CFI .95 and .93; RMSEA .087 and .108).

\subsubsection*{C.3.\ Scenario Order}

Participants saw the two scenarios in random order, and the order did not change the results. The split was even: 83 participants saw the less sensitive scenario first and 94 saw the more sensitive one first (binomial test, $p=.45$). The shift from the less to the more sensitive scenario did not differ by order for willingness, protection, effort, or benefit loss (Mann--Whitney tests, all $p_{\mathrm{Holm}}=1.00$). The share of participants who changed method between scenarios did not differ by order either (33 of 83 and 30 of 94; Fisher's exact test, $p=.35$).

\subsection*{D.\ Request and Reply Coding}

Two researchers independently coded the 138 sampled deletion requests on seven dimensions. Table~\ref{tab:coding} gives their agreement, measured with Gwet's AC1. Agreement was at least 0.89 on every dimension.

\begin{table}[t]
\centering
\caption{Agreement between the two coders on the 138 requests. Agreed is the number of requests on which both coders gave the same code. AC1 is Gwet's agreement coefficient.}
\label{tab:coding}
\small
\begin{tabular}{@{}p{0.40\columnwidth}rrr@{}}
\toprule
Dimension (what it codes) & Agreed & \% & AC1 \\
\midrule
Deletion location (where the request asks to delete: the conversation, memory, a backend database, the whole account, going forward, or unspecified) & 125/138 & 90.6 & 0.890 \\
Justification (whether the request gives a reason) & 131/138 & 94.9 & 0.943 \\
Data named (specific personal data, general personal data, or unspecified) & 133/138 & 96.4 & 0.956 \\
Mood (imperative, question, statement, or other) & 133/138 & 96.4 & 0.959 \\
Tone (polite, neutral, or friendly) & 137/138 & 99.3 & 0.990 \\
Verb (technical, lay, or no deletion verb) & 131/138 & 94.9 & 0.947 \\
Accompanying request (none, or asks alongside for verification, guidance on deleting it themselves, or keeping the context) & 133/138 & 96.4 & 0.962 \\
\midrule
Pooled (966 ratings) & 923/966 & 95.5 & 0.954 \\
\bottomrule
\end{tabular}
\end{table}

Each agent reply to a typed deletion request received one code, judged on what the agent said:
\begin{itemize}
\item complied: the reply confirms the request is done, with no stated limit and no redirect.
\item partial: the reply confirms an action but states a limit or redirects for the rest, or only promises not to use the fact.
\item deflected: the reply confirms no action and redirects, for example to Settings.
\item refused: the reply says it cannot or will not delete, with no redirect.
\item silently ignored: the reply engages with the topic but never acknowledges the request.
\item other: anything else, such as a clarifying question, a permission prompt, or a usage-limit screen.
\end{itemize}
The two coders agreed on 767 of 788 replies (Gwet's AC1 of 0.969).
